% Part: first-order-logic
% Chapter: beyond
% Section: intuitionistic-logic

\documentclass[../../../include/open-logic-section]{subfiles}

\begin{document}

\olfileid{fol}{byd}{il}

\olsection{منطق شهودگرایانه}

برخلاف منطق مرتبهٔ دوم و مراتب بالاتر، منطق مرتبهٔ اول شهودگرایانه
نسخه‌ای محدودشده از منطق کلاسیک است که برای بازنمایی گونه‌ای استدلال
«سازنده‌تر» در نظر گرفته شده است. مثال‌های زیر می‌توانند برخی از
انگیزه‌های بنیادین آن را روشن کنند.

فرض کنید روزی کسی نزد شما بیاید و اعلام کند عدد طبیعی~$x$ای را تعیین
کرده است که این خاصیت را دارد: اگر $x$ اول باشد، فرضیهٔ ریمان صادق است،
و اگر $x$ مرکب باشد، فرضیهٔ ریمان کاذب است. چه خبر خوبی!{} صادق یا کاذب
بودن فرضیهٔ ریمان یکی از مسائل بزرگ حل‌نشدهٔ ریاضیات است، و به نظر
می‌رسد آن شخص مسئله را به محاسبه، یعنی تعیین اول یا مرکب بودن عددی
معین، فروکاسته است.

مقدار جادویی~$x$ چیست؟ آن را چنین وصف می‌کند: $x$ عدد طبیعی‌ای است که
اگر فرضیهٔ ریمان صادق باشد برابر~$7$ است و در غیر این صورت برابر~$9$.

با خشم پولتان را پس می‌خواهید. از دیدگاه کلاسیک، توصیف بالا واقعاً
مقدار یکتایی برای~$x$ تعیین می‌کند؛ اما آنچه در حقیقت می‌خواهید مقداری
برای~$x$ است که emph{به‌صراحت} داده شده باشد.

مثالی دیگر و شاید کمتر ساختگی را در نظر بگیرید. می‌دانیم می‌توان عددی
گنگ را به توانی گویا رساند و نتیجه‌ای گویا به دست آورد. برای نمونه،
$\sqrt{2}^2 = 2$. آنچه کمتر روشن است این است که آیا می‌توان عددی گنگ
را به توانی emph{گنگ} رساند و نتیجه‌ای گویا به دست آورد یا نه. قضیهٔ
زیر پاسخی مثبت می‌دهد:

\begin{thm}
اعداد گنگ~$a$ و~$b$ای وجود دارند که $a^b$ گویاست.
\end{thm}

\begin{proof}
$\sqrt{2}^{\sqrt{2}}$ را در نظر بگیرید. اگر این عدد گویا باشد، کار
تمام است: می‌توانیم $a = b = \sqrt{2}$ بگیریم. در غیر این صورت، این
عدد گنگ است. آنگاه داریم
\[
(\sqrt{2}^{\sqrt{2}})^{\sqrt{2}} = \sqrt{2}^{\sqrt{2} \cdot
  \sqrt{2}} = \sqrt{2}^2 = 2,
\]
که بی‌تردید گویاست. پس در این حالت، $a$ را
$\sqrt{2}^{\sqrt{2}}$ و $b$ را~$\sqrt 2$ بگیرید.
\end{proof}

آیا این برهانی معتبر است؟ بیشتر ریاضی‌دانان آن را معتبر می‌دانند. اما
باز نکته‌ای تا اندازه‌ای ناخوشایند وجود دارد: وجود یک جفت عدد حقیقی با
خاصیتی معین را اثبات کرده‌ایم، بی‌آنکه بتوانیم بگوییم آن جفت
\emph{کدام} است. می‌توان همان نتیجه را چنان اثبات کرد که جفت~$a$ و~$b$
در خود برهان \emph{داده شود}: $a = \sqrt{3}$ و $b = \log_3 4$ بگیرید.
آنگاه
\[
a^b = \sqrt{3}^{\log_3 4} = 3^{1/2 \cdot \log_3 4} = (3^{\log_3
  4})^{1/2} = 4^{1/2}= 2,
\]
زیرا $3^{\log_3 x} = x$. هر دو عدد گنگ‌اند: گنگ‌بودن~$\sqrt3$ معلوم
است، و اگر $\log_3 4=p/q$ برای اعداد صحیح مثبت~$p$ و~$q$ بود، از آن
$3^p=4^q=2^{2q}$ نتیجه می‌شد که با یکتایی تجزیه به عوامل اول ناسازگار
است.

منطق شهودگرایانه برای بازنمایی گونه‌ای استدلال طراحی شده است که حرکت‌هایی
مانند حرکت برهان نخست را مجاز نمی‌شمارد. اثبات وجود~$x$ای که~$!A(x)$
را برآورده می‌کند بدین معناست که باید~$x$ معینی و برهانی بدهید که نشان
دهد آن~$!A$ را برآورده می‌کند، چنان‌که در برهان دوم چنین بود. اثبات
اینکه $!A$ یا $!B$ برقرار است مستلزم آن است که بتوانید یکی از آن دو را
اثبات کنید.

به بیان صوری، منطق مرتبهٔ اول شهودگرایانه چیزی است که با محدودکردن
دستگاه !!a{derivation} منطق مرتبهٔ اول به شیوه‌ای معین به دست می‌آید.
به همین ترتیب، نسخه‌هایی شهودگرایانه از منطق مرتبهٔ دوم و مراتب بالاتر
وجود دارند. از دیدگاه ریاضی، این‌ها صرفاً دستگاه‌های استنتاجی صوری‌اند،
اما چنان‌که گفته شد برای بازنمایی گونه‌ای استدلال ریاضی در نظر گرفته
شده‌اند. می‌توان این استدلال را همان گونه‌ای دانست که دیدگاهی فلسفی
دربارهٔ ریاضیات (مانند شهودگرایی براور) توجیه می‌کند؛ می‌توان آن را
گونه‌ای استدلال ریاضی «انضمامی‌تر» و رضایت‌بخش‌تر دانست (در امتداد
ساخت‌گرایی بیشاپ)؛ و می‌توان دربارهٔ اینکه توصیف صوری تا چه اندازه
انگیزهٔ غیرصوری را به‌درستی بازمی‌نماید بحث کرد. اما هر موضع فلسفی‌ای
که داشته باشیم، می‌توانیم منطق شهودگرایانه را همچون منطقی با عرضهٔ صوری
بررسی کنیم؛ و به هر دلیلی، بسیاری از منطق‌دانان ریاضی این بررسی را
جالب می‌یابند.

تفسیری سازنده و غیرصوری از رابط‌های شهودگرایانه وجود دارد که معمولاً
تفسیر BHK نامیده می‌شود (برگرفته از نام‌های براور، هایتینگ و
کولموگروف). این تفسیر چنین است: برهان $!A \land !B$ از برهان~$!A$ همراه
با برهان~$!B$ تشکیل می‌شود؛ برهان $!A \lor !B$ یا برهان~$!A$ است یا
برهان~$!B$، و اطلاعات صریحی داریم که کدام حالت برقرار است؛ برهان
$!A \lif !B$ رویه‌ای است که برهان~$!A$ را به برهان~$!B$ تبدیل می‌کند؛
برهان $\lforall[x][!A(x)]$ رویه‌ای است که برای هر مقدار~$x$ برهانی برای
$!A(x)$ بازمی‌گرداند؛ و برهان $\lexists[x][!A(x)]$ از مقداری برای~$x$
همراه با برهانی تشکیل می‌شود که نشان می‌دهد آن مقدار~$!A$ را برآورده
می‌کند. می‌توان این تفسیر را با اصطلاحات محاسباتی‌ای وصف کرد که به
«هم‌ریختی کری--هاوارد» یا «انگارهٔ !!{formula}s همچون نوع‌ها» معروف
است: !!a{formula} را مشخص‌کنندهٔ گونه‌ای نوع داده و برهان‌ها را اشیای
محاسباتی از این نوع‌های داده در نظر بگیرید که به ما امکان می‌دهند صدق
!!{formula} متناظر را ببینیم.

منطق شهودگرایانه را غالباً منطق کلاسیک «منهای» اصل طرد شق ثالث
می‌دانند. قضیهٔ زیر این مطلب را دقیق‌تر می‌کند.
\begin{thm}
طرح‌واره‌های اصل موضوع زیر از دیدگاه شهودگرایانه هم‌ارزند:
\begin{enumerate}
\item $(\lnot !A \lif \lfalse) \lif !A$.
\item $!A \lor \lnot !A$
\item $\lnot \lnot !A \lif !A$
\end{enumerate}
\end{thm}
به‌دست‌آوردن نمونه‌های هر یک از این طرح‌واره‌ها از هر یک از دو طرح‌وارهٔ
دیگر تمرین خوبی در منطق شهودگرایانه است.

نخستین دستگاه‌های استنتاجی منطق گزاره‌ای شهودگرایانه، که به‌عنوان
صورت‌بندی‌های صوری شهودگرایی براور عرضه شدند، مستقلاً کار کولموگروف،
گلیونکو و هایتینگ‌اند. نخستین صورت‌بندی صوری منطق مرتبهٔ اول
شهودگرایانه (و بخش‌هایی از ریاضیات شهودگرایانه) از آنِ هایتینگ است.
هرچند شماری از طرح‌واره‌های معتبر کلاسیک از دیدگاه شهودگرایانه معتبر
نیستند، بسیاری از آن‌ها معتبرند.

\emph{ترجمهٔ نقیض مضاعف} رابطه‌ای مهم میان منطق کلاسیک و شهودگرایانه
را توصیف می‌کند. این ترجمه به‌صورت استقرایی چنین تعریف می‌شود
($!A^N$ را ترجمهٔ «شهودگرایانه»ی !!{formula} کلاسیک~$!A$ در نظر
بگیرید):
\begin{align*}
!A^N & \ident \lnot\lnot !A \quad \text{برای !!{formula}s اتمی $!A$} \\
(!A \land !B)^N & \ident (!A^N \land !B^N) \\
(!A \lor !B)^N & \ident  \lnot\lnot (!A^N \lor !B^N) \\
(!A \lif !B)^N & \ident (!A^N \lif !B^N) \\
(\lforall[x][!A])^N & \ident \lforall[x][!A^N] \\
(\lexists[x][!A])^N & \ident \lnot\lnot\lexists[x][!A^N]
\end{align*}
کولموگروف و گلیونکو صورت‌هایی از این ترجمه را برای منطق گزاره‌ای عرضه
کردند؛ صورت منطق محمولات مستقلاً از آنِ گودل و گنتسن است. داریم:

\begin{thm}
\begin{enumerate}
\item $!A \liff !A^N$ به‌صورت کلاسیک اثبات‌پذیر است.
\item اگر $!A$ به‌صورت کلاسیک اثبات‌پذیر باشد، آنگاه $!A^N$ به‌صورت
  شهودگرایانه اثبات‌پذیر است.
\end{enumerate}
\end{thm}

اکنون می‌توانیم گفت‌وگوی زیر را مجسم کنیم. ریاضی‌دان کلاسیک: «$!A$ را
اثبات کرده‌ام!» ریاضی‌دان شهودگرا: «برهانت معتبر نیست. آنچه واقعاً
اثبات کرده‌ای $!A^N$ است.» ریاضی‌دان کلاسیک: «برای من فرقی ندارد!» از
نظر ریاضی‌دان کلاسیک، شهودگرا صرفاً مو را از ماست می‌کشد، زیرا آن دو
هم‌ارزند. اما شهودگرا اصرار دارد که تفاوتی وجود دارد.

توجه کنید که ترجمهٔ بالا تنها به منطق محض مربوط است؛ این ترجمه به این
پرسش نمی‌پردازد که اصول موضوع emph{غیرمنطقی} مناسب برای ریاضیات کلاسیک
و شهودگرایانه کدام‌اند یا چه رابطه‌ای میان آن‌ها برقرار است. اما بسط
اندک زیر از قضیهٔ بالا اطلاعات سودمندی به دست می‌دهد:

\begin{thm}
اگر $\Gamma$، $!A$ را به‌صورت کلاسیک اثبات کند، $\Gamma^N$، $!A^N$ را
به‌صورت شهودگرایانه اثبات می‌کند.
\end{thm}

به بیان دیگر، اگر $!A$ از فرض‌هایی به‌صورت کلاسیک اثبات‌پذیر باشد،
$!A^N$ از ترجمه‌های نقیض مضاعف آن فرض‌ها اثبات‌پذیر است.

برای نشان‌دادن اینکه یک جمله یا !!{formula} گزاره‌ای از دیدگاه
شهودگرایانه معتبر است، تنها کافی است برهانی ارائه کنید. اما چگونه
می‌توانید نشان دهید که معتبر نیست؟ برای این منظور به معناشناسی‌ای نیاز
داریم که صحیح و، ترجیحاً، تام باشد. معناشناسی‌ای از کریپکی به‌خوبی این
نیاز را برآورده می‌کند.

می‌توانیم همان برنامه‌ای را اجرا کنیم که برای منطق کلاسیک اجرا کردیم:
معناشناسی را تعریف کنیم و صحت و تمامیت را اثبات کنیم. بااین‌حال، توجه
به تمایز زیر ارزشمند است. در منطق کلاسیک، معناشناسی به معنایی که در
معنای رابط‌ها نهفته بود، معناشناسی «بدیهی» به شمار می‌آمد. هرچند
می‌توان برای معناشناسی کریپکی انگیزه‌ای شهودی فراهم کرد، این معناشناسی
همان احساس ناگزیربودن را پدید نمی‌آورد. افزون بر این، مفهوم
!!{structure} کلاسیک مفهومی طبیعی در ریاضیات است؛ پس می‌توانیم مفهوم
!!a{structure} را ابزاری برای مطالعهٔ منطق مرتبهٔ اول کلاسیک بدانیم، یا
منطق مرتبهٔ اول کلاسیک را ابزاری برای مطالعهٔ !!{structure}s ریاضی.
در مقابل، !!{structure}s کریپکی را تنها می‌توان سازه‌هایی منطقی دانست؛
به نظر نمی‌رسد آن‌ها مستقلاً مورد علاقهٔ ریاضی باشند.

یک !!{structure} کریپکی~$\mModel M = \tuple{W, R, V}$ برای زبانی
گزاره‌ای از مجموعه‌ای چون~$W$، ترتیب جزئی~$R$ روی~$W$ با کمترین
!!{element}، و گمارشی «یکنوا» از متغیرهای گزاره‌ای به !!{element}s
مجموعهٔ~$W$ تشکیل می‌شود. شهود این است که !!{element}s مجموعهٔ~$W$
نمایندهٔ «جهان‌ها» یا «حالت‌های معرفت»اند؛ عنصر $v \geq u$ نمایندهٔ
«حالت ممکن آینده»ی~$u$ است؛ و متغیرهای گزاره‌ای گمارده‌شده به~$u$
گزاره‌هایی هستند که در حالت~$u$ صادق دانسته می‌شوند. رابطهٔ اجبار
$\mSat{M}{!A}[w]$ سپس این رابطه را به !!{formula}s دلخواه زبان گسترش
می‌دهد؛ $\mSat{M}{!A}[w]$ را چنین بخوانید: «$!A$ در حالت~$w$ صادق
است.» این رابطه به‌صورت استقرایی چنین تعریف می‌شود:
\begin{enumerate}
\item $\mSat{M}{\Obj p_i}[w]$ هرگاه و تنها هرگاه $\Obj p_i$ یکی از
  متغیرهای گزاره‌ای گمارده‌شده به~$w$ باشد.
\item $\mSat/{M}{\lfalse}[w]$.
\item $\mSat{M}{(!A \land !B)}[w]$ هرگاه و تنها هرگاه
  $\mSat{M}{!A}[w]$ و $\mSat{M}{!B}[w]$.
\item $\mSat{M}{(!A \lor !B)}[w]$ هرگاه و تنها هرگاه
  $\mSat{M}{!A}[w]$ یا $\mSat{M}{!B}[w]$.
\item $\mSat{M}{(!A \lif !B)}[w]$ هرگاه و تنها هرگاه، برای هر
  $w' \geq w$، اگر $\mSat{M}{!A}[w']$، آنگاه $\mSat{M}{!B}[w']$.
\end{enumerate}
این تمرین خوبی است که بکوشید با ساختن !!{structure} کریپکی‌ای که مثال
نقضی به دست می‌دهد نشان دهید $\lnot (p \land q) \lif
(\lnot p \lor \lnot q)$ از دیدگاه شهودگرایانه معتبر نیست.

\end{document}
